\section{Introduction}

A real Hadamard matrix of order \(n\) is a matrix \(H\) with entries
in \(\{-1,1\}\) such that \(HH^{\mathsf T}=nI_n\). It is
\emph{circulant} if there are signs \(h_0,\ldots,h_{n-1}\) for which
\[
 H_{ij}=h_{j-i\bmod n}\qquad(0\leq i,j<n).
\]
Thus every row is a cyclic shift of the first. For such a sign matrix,
define the periodic autocorrelations by
\[
 P_h(t)=\sum_{j=0}^{n-1}h_jh_{j+t\bmod n}
 \qquad(0\leq t<n).
\]
Here \(P_h(0)=n\), and the Hadamard condition is equivalent to
\(P_h(t)=0\) for every \(1\leq t<n\).
The circulant Hadamard conjecture, traditionally attributed to Ryser
\cite{Ryser1963,LeungSchmidt2012}, asserts that the only possible
orders are \(1\) and \(4\). We prove this conjecture.

\begin{theorem}\label{thm:main}
For a positive integer \(n\), a real circulant Hadamard matrix of
order \(n\) exists if and only if \(n\in\{1,4\}\).
\end{theorem}

The theorem also completes the classification of Barker-sequence
lengths. For an integer \(n>1\), a sign sequence
\(a=(a_0,\ldots,a_{n-1})\in\{-1,1\}^n\) is a \emph{Barker sequence}
if its nontrivial aperiodic autocorrelations satisfy
\[
 |C_a(t)|\leq1\quad(1\leq t<n),\qquad
 C_a(t)=\sum_{j=0}^{n-t-1}a_ja_{j+t}.
\]
The Barker-sequence conjecture asserts that no such sequence has
length greater than \(13\). Turyn and Storer established the odd-length
nonexistence theorem \cite{TurynStorer1961}; we use the later proof of
Schmidt and Willms \cite{SchmidtWillms2016} for the exact list.
For even length greater
than two, the Barker inequalities imply that all nontrivial periodic
autocorrelations vanish \cite{TurynStorer1961,Turyn1965}.
Theorem~\ref{thm:main} therefore settles the remaining even case.

\begin{corollary}[Barker-sequence lengths]\label{cor:barker}
For an integer \(n>1\), a Barker sequence of length \(n\) exists if
and only if
\[
 n\in\{2,3,4,5,7,11,13\}.
\]
\end{corollary}
